% =============================================================================
% ЛЕКЦИЯ 2 - Примеры кода
% =============================================================================

\lecture{2}{Примеры кода на разных языках}

\lecturedate{8 сентября 2025}

Эта лекция демонстрирует возможности подсветки синтаксиса для различных языков программирования.

\section{Python}

Python — высокоуровневый язык программирования с динамической типизацией.

\begin{python}
from typing import List, Optional
from dataclasses import dataclass

@dataclass
class Point:
    x: float
    y: float

    def distance(self, other: "Point") -> float:
        """Вычисляет расстояние до другой точки."""
        return ((self.x - other.x)**2 + (self.y - other.y)**2) ** 0.5

def find_closest(points: List[Point], target: Point) -> Optional[Point]:
    if not points:
        return None
    return min(points, key=lambda p: p.distance(target))

# Пример использования
points = [Point(0, 0), Point(3, 4), Point(1, 1)]
target = Point(2, 2)
closest = find_closest(points, target)
print(f"Ближайшая точка: {closest}")  # Point(x=1, y=1)
\end{python}

\section{C++}

Современный C++ (C++17/20) с использованием STL и умных указателей.

\begin{cpp}
#include <iostream>
#include <vector>
#include <memory>
#include <algorithm>

template<typename T>
class Stack {
private:
    std::vector<T> data;

public:
    void push(T value) {
        data.push_back(std::move(value));
    }

    std::optional<T> pop() {
        if (data.empty()) return std::nullopt;
        T value = std::move(data.back());
        data.pop_back();
        return value;
    }

    [[nodiscard]] bool empty() const noexcept {
        return data.empty();
    }
};

int main() {
    auto stack = std::make_unique<Stack<int>>();
    stack->push(42);

    if (auto val = stack->pop()) {
        std::cout << "Popped: " << *val << std::endl;
    }
    return 0;
}
\end{cpp}

\section{Rust}

Rust — язык системного программирования с гарантиями безопасности памяти.

\begin{rust}
use std::collections::HashMap;

#[derive(Debug, Clone)]
struct User {
    name: String,
    age: u32,
}

impl User {
    fn new(name: impl Into<String>, age: u32) -> Self {
        Self { name: name.into(), age }
    }
}

fn group_by_age(users: &[User]) -> HashMap<u32, Vec<&User>> {
    let mut groups: HashMap<u32, Vec<&User>> = HashMap::new();

    for user in users {
        groups.entry(user.age).or_default().push(user);
    }

    groups
}

fn main() {
    let users = vec![
        User::new("Alice", 25),
        User::new("Bob", 30),
        User::new("Charlie", 25),
    ];

    let grouped = group_by_age(&users);
    println!("{:#?}", grouped);
}
\end{rust}

\section{Go}

Go — компилируемый язык от Google с встроенной поддержкой конкурентности.

\begin{golang}
package main

import (
    "fmt"
    "sync"
)

type Counter struct {
    mu    sync.Mutex
    value int
}

func (c *Counter) Increment() {
    c.mu.Lock()
    defer c.mu.Unlock()
    c.value++
}

func (c *Counter) Value() int {
    c.mu.Lock()
    defer c.mu.Unlock()
    return c.value
}

func main() {
    counter := &Counter{}
    var wg sync.WaitGroup

    for i := 0; i < 1000; i++ {
        wg.Add(1)
        go func() {
            defer wg.Done()
            counter.Increment()
        }()
    }

    wg.Wait()
    fmt.Printf("Final value: %d\n", counter.Value())
}
\end{golang}

\section{JavaScript/TypeScript}

Пример асинхронного кода на JavaScript.

\begin{javascript}
interface User {
    id: number;
    name: string;
    email: string;
}

async function fetchUsers(): Promise<User[]> {
    const response = await fetch('/api/users');
    if (!response.ok) {
        throw new Error(`HTTP error: ${response.status}`);
    }
    return response.json();
}

const processUsers = async () => {
    try {
        const users = await fetchUsers();
        const names = users.map(u => u.name);
        console.log('Users:', names.join(', '));
    } catch (error) {
        console.error('Failed to fetch users:', error);
    }
};

// Использование Promise.all для параллельных запросов
const [users, posts] = await Promise.all([
    fetchUsers(),
    fetchPosts()
]);
\end{javascript}

\section{SQL}

Пример сложного SQL-запроса с CTE и оконными функциями.

\begin{sql}
WITH monthly_sales AS (
    SELECT
        DATE_TRUNC('month', order_date) AS month,
        product_id,
        SUM(quantity * price) AS revenue
    FROM orders
    WHERE order_date >= '2024-01-01'
    GROUP BY 1, 2
),
ranked_products AS (
    SELECT
        month,
        product_id,
        revenue,
        ROW_NUMBER() OVER (
            PARTITION BY month
            ORDER BY revenue DESC
        ) AS rank
    FROM monthly_sales
)
SELECT
    r.month,
    p.name AS product_name,
    r.revenue,
    r.rank
FROM ranked_products r
JOIN products p ON r.product_id = p.id
WHERE r.rank <= 5
ORDER BY r.month, r.rank;
\end{sql}

\section{Haskell}

Функциональное программирование на Haskell.

\begin{haskell}
module Main where

import Data.List (sortBy)
import Data.Ord (comparing)

-- Определение типа данных
data Tree a = Empty | Node a (Tree a) (Tree a)
    deriving (Show, Eq)

-- Вставка в бинарное дерево поиска
insert :: Ord a => a -> Tree a -> Tree a
insert x Empty = Node x Empty Empty
insert x (Node y left right)
    | x < y     = Node y (insert x left) right
    | otherwise = Node y left (insert x right)

-- Обход дерева (in-order)
inorder :: Tree a -> [a]
inorder Empty = []
inorder (Node x left right) = inorder left ++ [x] ++ inorder right

-- Создание дерева из списка
fromList :: Ord a => [a] -> Tree a
fromList = foldr insert Empty

main :: IO ()
main = do
    let tree = fromList [5, 3, 7, 1, 4, 6, 8]
    print $ inorder tree  -- [1,3,4,5,6,7,8]
\end{haskell}

\section{Bash}

Пример shell-скрипта для автоматизации.

\begin{bash}
#!/bin/bash
set -euo pipefail

# Функция для логирования
log() {
    echo "[$(date '+%Y-%m-%d %H:%M:%S')] $*"
}

# Проверка зависимостей
check_deps() {
    local deps=("git" "docker" "jq")
    for dep in "${deps[@]}"; do
        if ! command -v "$dep" &> /dev/null; then
            log "ERROR: $dep is not installed"
            exit 1
        fi
    done
}

# Основная логика
main() {
    check_deps

    log "Starting deployment..."

    # Получение последней версии
    VERSION=$(git describe --tags --always)
    log "Version: $VERSION"

    # Сборка и деплой
    docker build -t "myapp:$VERSION" .
    docker push "myapp:$VERSION"

    log "Deployment complete!"
}

main "$@"
\end{bash}
